\section{从函数到层}

\ob{3.1}{ob:3.1}
下一步，把 $\Bun_n$ 不只看成集合，而看成定义在 $\Fq$ 上代数簇的有理点集，并考虑其上的 $\ell$-进层（$\ell$-进层的简述见 \S4）。

一个小麻烦是：$\Bun_n$ 背后的几何对象不是簇，而是代数叠（参看 [L-MB]），记为 $\sBun_n$。叠与概形的主要差别如下。

按 Yoneda，范畴的对象（此处即概形 $\mathcal{Y}$）由函子 $\mathrm{Schemes}\Rightarrow\mathrm{Sets}$，$S\mapsto\Hom(S,\mathcal{Y})$ 完全决定。

在 $\sBun_n$ 的情形，上述 Hom 确实有定义，但不再是集合，而是范畴。这一点在讨论 $\Bun_n$ 时已经可见：考虑 $X$ 上向量丛的范畴，比考虑向量丛同构类的集合自然得多。

一般地，对概形 $S$，定义 $\Hom(S,\sBun_n)$ 为范畴，其对象是 $X\times S$ 上秩 $n$ 的丛，态射是向量丛之间的同构。这样定义的 $\sBun_n$ 确实是代数叠。

然而，从 $\ell$-进层的观点看，叠的行为如同代数簇。特别地，若 $\mathcal{Y}$ 是叠，范畴 $D^b(\mathcal{Y})$ 有意义（参看 \S4）。此外，若 $Y$ 是范畴 $\Hom(\Spec(\Fq),\mathcal{Y})$ 中对象的同构类集合，则 $D^b(\mathcal{Y})$ 的每个对象 $\Fc$ 都定义 $Y$ 上的函数 $f$，且层-函数对应的全部性质都成立。

特别地，范畴 $D^b(\sBun_n)$ 有意义，且每个对象 $\Fc\in D^b(\sBun_n)$ 都给出 $\Bun_n$ 上的函数 $f$。

\begin{annote}
若只记住「叠是取值到群胚的函子」，上面的说法会显得突兀：既然 $\Hom(S,\sBun_n)$ 是范畴而非集合，为什么还能谈 $D^b(\sBun_n)$？关键在于 $\ell$-进层只看见叠的「局部」信息：叠有光滑覆盖，而层的拉回沿光滑覆盖是下降的，于是可以像在簇上一样定义可构造层及其导出范畴。代价是：$D^b$ 中的对象不再有「逐点取值」的意义，只有沿 $\Fq$-点上的纤维取迹才得到函数。\par\vspace{0.15em}{\raggedleft 【来源：Laumon--Moret-Bailly, \emph{Champs alg\'ebriques}（A）；叠上 \'etale 层的一般约定（B）】\par}
\end{annote}

\ob{3.2}{ob:3.2}
现在讨论 Hecke 算子在这个几何情境中变成什么。

先固定 $x\in|X|$。考虑 Hecke 叠 $\sH_x^i$：$\Hom(S,\sH_x^i)$ 是三元组 $(M,M',\beta)$ 的范畴，其中 $M$ 与 $M'$ 是 $X\times S$ 上秩 $n$ 向量丛，$\beta$ 是凝聚层的嵌入 $M\hookrightarrow M'$，使得 $M'/M$ 概形论地支撑在 $x\times S\subset X\times S$ 上，且是该子概形上秩 $i$ 的向量丛。

$\sH_x^i$ 的 $\Fq$-点的同构类集合，就是我们原来的 $H_x^i$。和前面一样，有映射 $h^{\leftarrow},h^{\rightarrow}:\sH_x^i\to\sBun_n$。但现在对 $\sBun_n$ 的给定点，即向量丛 $M$，原像 $(h^{\leftarrow})^{-1}\subset\sH_x^i$ 不再是集合，而是同构于 $\Gr^i(M_x)$ 的概形。

引入 Hecke 函子 $T_x^i:D^b(\sBun_n)\to D^b(\sBun_n)$：
\[
\Fc\mapsto(h^{\rightarrow})_!((h^{\leftarrow})^*(\Fc))[i(n-i)],
\]
其中括号内的表达式表示上同调平移。

\ob{3.3}{ob:3.3}
不过，让点 $x$ 沿曲线移动会更好。引入「大」Hecke 叠 $\sH^i$ 如下。

$\Hom(S,\sH^i)$ 是四元组 $(\phi,M,M',\beta)$ 的范畴，其中 $M,M'$ 同前，$\phi$ 是映射 $S\to X$，$\beta$ 仍是凝聚层的嵌入 $M\hookrightarrow M'$，使得商支撑在图像 $\Gamma_\phi\subset X\times S$ 上，且是其上秩 $i$ 的向量丛。

除原有的投影 $h^{\leftarrow},h^{\rightarrow}:\sH^i\to\sBun_n$ 外，现在还有自然映射 $s:\sH^i\to X$，把上述四元组 $(\phi,M,M',\beta)$ 送到 $\phi$。对固定 $x\in X$，有 $s^{-1}(x)\simeq\sH_x^i$。

定义 Hecke 函子 $T^i:D^b(\sBun_n)\to D^b(X\times\sBun_n)$ 为
\[
\Fc\mapsto(s\times h^{\rightarrow})_!((h^{\leftarrow})^*(\Fc))[i(n-i)].
\]

\ob{3.4}{ob:3.4}
现在可以给出 Hecke 本征层的定义。我们的定义稍嫌粗糙（对任意群都适用的更好定义见 [BD]）。

设 $\sigma:\pi_1(X)\to\GL_n(\Ql)$ 是表示；回忆 $E_\sigma$ 是相应的 $X$ 上 $n$ 维局部系统。

若对每个 $i=1,\dots,n$ 都存在同构
\[
T^i(\Fc_\sigma)\simeq\Lambda^i(E_\sigma)\boxtimes\Fc_\sigma,
\]
则称对象 $\Fc_\sigma\in D^b(\sBun_n)$ 为关于 $\sigma$ 的 Hecke 本征层；其中 $\Lambda^i$ 表示 $i$ 次外幂。

若基域 $k=\Fq$，且 $\Fc_\sigma$ 是关于 $\sigma$ 的 Hecke 本征层，则其对应的 $\GL_n(\cK)\backslash\GL_n(\A)/\GL_n(\sO)$ 上的函数，是 Hecke 算子的公共特征函数，特征值由 $\sigma$ 给出，如 \S1.9 所述。

几何朗兰兹猜想（由 [L] 在 Drinfeld 原始论文之后提出；[Dr] 完整处理了 $\GL_2$ 的情形）预测：对每个如上所述的不可约 $\sigma$，确实存在 Hecke 本征层 $\Fc_\sigma$，而且它是尖点的（即 \S2.8 所定义的函子 $r^n_{n_1,n_2}$ 在 $\Fc_\sigma$ 上为零）。

\begin{annote}
原书自称上述定义「稍嫌粗糙」，并在 [BD] 中指出对任意群都适用的更好定义。差别大致在于：这里只对 $\GL_n$ 陈述，且把本征性直接写成外幂 $\Lambda^i(E_\sigma)$ 的显式同构。对一般约化群 $G$，Hecke 函子作用在 $\Bun_G$ 上，本征层的「特征值」应取对偶群 $\check{G}$ 的局部系统，而外幂 $\Lambda^i$ 只是 $G=\GL_n$ 时标准表示的外幂。\par\vspace{0.15em}{\raggedleft 【来源：原书 \S3.4 与 [BD]（A）】\par}
\end{annote}

\ob{3.5}{ob:3.5}
最后讨论一个问题：为什么把讨论限制在球（即无分歧）情形？

理由如下：如前所述，球尖点自守表示对应 $\GL_n(\cK)\backslash\GL_n(\A)/\GL_n(\sO)$ 上一个确定的函数 $f$，它唯一到相差一个标量。因此，几何上我们要构造 $D^b(\sBun_n)$ 中一个特定的对象；当 $k=\Fq$ 时，它给出这个函数。

对一般的尖点自守表示 $\pi$，很难在 $\pi\subset\Funct_{\cusp}(\GL_n(\cK)\backslash\GL_n(\A))$ 中典范地选出一个函数。\footnote{由于我们的群是 $\GL_n$，这样的选择存在，即所谓的新形式，但从几何观点看并不那么自然。}因此，在相应的 $\sBun_n$ 版本上该构造什么样的层，并不清楚。

\section{附录：层-函数对应}

\ob{4.1}{ob:4.1}
这里回顾 Grothendieck 的层-函数对应的若干事实。

首先，若 $\mathcal{Y}$ 是（任意基域上的）概形，可以定义范畴 $D^b(\mathcal{Y})$，称为上同调可构造的 $\ell$-进层的导出范畴。

若 $\phi:\mathcal{Y}_1\to\mathcal{Y}_2$ 是概形间的态射，则有拉回与直接像函子
\[
\phi^*:D^b(\mathcal{Y}_2)\to D^b(\mathcal{Y}_1),\qquad
\phi_!:D^b(\mathcal{Y}_1)\to D^b(\mathcal{Y}_2).
\]

若 $\Fc_1$ 与 $\Fc_2$ 是 $D^b(\mathcal{Y})$ 的两个对象，可以定义它们的张量积 $\Fc_1\otimes\Fc_2\in D^b(\mathcal{Y})$。

\ob{4.2}{ob:4.2}
现设 $\mathcal{Y}$ 定义在基域 $\Fq$ 上。记 $Y$ 为 $\mathcal{Y}$ 的 $\Fq$-值点集（它是有限集），$\Funct(Y)$ 为其上 $\Ql$-值函数构成的向量空间。$D^b(\mathcal{Y})$ 的每个对象 $\Fc$ 都给出 $\Funct(Y)$ 中一个元素 $f$。

首先，若 $\mathcal{Y}$ 是点概形 $\Spec(\Fq)$，则 $D^b(\Spec(\Fq))$ 的对象就是有限维 $\Ql$-向量空间构成的有复形，带 $\mathrm{Gal}(\Flq/\Fq)$ 的作用。对应的函数（此时只是 $\Ql$ 的一个元素）是 $\mathrm{Gal}(\Flq/\Fq)$ 中 Frobenius 元的迹的交替和。对一般的 $\mathcal{Y}$、$\Fc\in D^b(\mathcal{Y})$ 与 $y\in Y$，$f$ 在 $y$ 处的值等于上述规则赋给 $\Fc_y$ 的数，其中 $\Fc_y\in D^b(\Spec(\Fq))$ 是 $\Fc$ 沿对应于 $y$ 的态射 $\Spec(\Fq)\to\mathcal{Y}$ 的拉回。

对应 $\Fc\mapsto f$ 的基本性质如下。张量积 $\Fc_1\otimes\Fc_2$ 对应 $f_1\cdot f_2$（逐点乘积）。对映射 $\phi:\mathcal{Y}_1\to\mathcal{Y}_2$，拉回 $\phi^*(\Fc)$ 对应 $\phi^*(f)$，推前同理。（就证明而言，前两条性质由定义立得，第三条即 Grothendieck--Lefschetz 迹公式。）

\begin{annote}
第三条性质值得记住：它把「沿态射推前」与「逐纤维求和」等同起来。前文 \S2.5 中 $(h^{\rightarrow})_!$ 的求和、\S2.8 中 $q_!$ 的积分，都是它在函数侧的表现；到了层侧，同一件事由上同调的直接像承担，而迹公式保证两边算出的函数一致。\par\vspace{0.15em}{\raggedleft 【来源：原书 \S4.2（A）；Grothendieck--Lefschetz 迹公式的标准表述（B）】\par}
\end{annote}

\ob{4.3}{ob:4.3}
事实证明，许多情形下定义在 $\mathcal{Y}/\Fq$ 的 $Y$ 上的「有趣」函数，都按上述方式来自 $\ell$-进层。某些情形下也可以反向构造。下面是一个例子。

设 $A$ 是连通交换代数群。若对象 $\mathcal{X}\in D^b(A)$ 是一维局部系统，其在单位元处的纤维平凡，且 $m^*(\mathcal{X})\simeq\mathcal{X}\boxtimes\mathcal{X}$（其中 $m:A\times A\to A$ 是乘法映射），则称 $\mathcal{X}$ 为特征层。

此时，对应的有限群 $A=A(\Fq)$ 上的函数 $\chi$ 是特征。

反之，给定 $A$ 的特征 $\chi$，可如下重构 $\mathcal{X}$。考虑 Lang 同态 $L:A\to A$，
\[
a\mapsto\mathrm{Fr}(a)/a,
\]
其中 $\mathrm{Fr}:A\to A$ 是 Frobenius 自同态。则 $L$ 是 \'etale Galois 覆盖，其覆盖变换群等于 $\ker(L)=A$。

所求的特征层 $\mathcal{X}$，由 $\pi_1(A)$ 的如下一维表示所对应的 $\ell$-进层给出：
\[
\pi_1(A)\to\ker(L)=A\xrightarrow{\chi}\Ql^\times .
\]

把这一构造用于 $A$ 是仿射直线 $\mathbb{A}^1$ 的情形，得到 $\mathbb{A}^1$ 上称为 Artin--Schreier 层的 $\ell$-进层。类似地，$A$ 是乘法群 $\mathbb{G}_m$ 时得到其上的 Kummer 层。

另一个例子是 $A=\Pic^0(X)$：由此可恢复 Lang 对函数域上整体类域论的证明中所用的构造；参看 [Se]。
